A further experimental reading concerns the relationship between precision and operational trust in the BN layer. Precision near $0.806$ at FPR near $0.027$ means that when the BN alerts under budget, the alert is usually associated with a truly severe project. That property can matter for escalation culture even when recall is only $0.450$: analysts may trust BN alerts more than a high-recall low-precision rule, provided they understand that silence is not a clean bill of health~\cite{Vickers2006,Saito2015}.

The champion's specificity of $0.964$ plays an analogous role for early warning. Most non-severe projects remain quiet at $\tau{=}0.51$, which limits alert fatigue at portfolio scale. Combined with decision-curve dominance over alert-all, the operating point is compatible with routine triage rather than with emergency-only use~\cite{Vickers2006,Gneiting2007}.

Feature provenance for the champion deserves the same scrutiny as BN parameters. Prefix variance statistics are computed only from early updates; agency/phase/category indicators are observed attributes; budget covariates come from the companion open dataset. No TEST-derived statistic enters feature construction. That statement is trivial in principle and frequently violated in practice when practitioners standardize on the full sample or select features with TEST feedback~\cite{Egwim2021,Tagharobi2025}.

The NLP fallback's weak macro-F1 is not hidden in an appendix because it is scientifically informative. It shows that unconstrained text classification on scarce unique reason strings does not create reliable BN parents. The regex-first design is therefore part of the method's validity argument, not merely an implementation convenience~\cite{Onisko2001,Woldesenbet2016}.

Bootstrap CI width differences across metrics also carry interpretation. Accuracy and ROC-AUC are stable; precision and recall move more under resampling because they depend on the minority class. Any deployment dashboard that shows a single precision point without uncertainty is overconfident relative to the evidence in Table~\ref{tab:boot}~\cite{Hand2009,GneitingRaftery2007}.

Policy AUC values near $0.872$ for BN posteriors are respectable for a sparse evidence model but remain well below the champion's $0.981$. That gap is expected: the champion sees continuous variance dynamics; the BN sees binary factor evidence under budget. The gap should discourage replacing early warning with BN inspection, and should encourage using them jointly~\cite{Luu2009,Egwim2021}.

Myopic VoI's limitation is real even when the scoreboard ties. In denser evidence regimes, myopic selection can be suboptimal relative to multi-step plans that reserve budget for later decisive inspections~\cite{Memarzadeh2016,Fauriat2020}. The present study cannot quantify that suboptimality on NYC data because even myopic and fixed policies already coincide. Multi-step VoI is therefore a high-value extension for future denser datasets, not a missing fix for the current tie.

Absence clamping as a missingness model is similarly provisional. Treating never-mentioned factors as absent is faithful to a hard reading of administrative silence and keeps inference simple. A softer reading would place a prior on silence and update it when inspections occur. Soft silence models could increase policy differentiation by enlarging effective uncertainty, at the cost of additional parameters that open data may not identify~\cite{Heckerman1996,Song2021}.

The eight-factor ontology's compression of survey taxonomies is a modeling decision with trade-offs. Finer ontologies could improve explainability for specialists and could match contractual clause libraries. They would also increase the fraction of never-mentioned factors and could worsen sparsity. Until reason text becomes richer, compression is a rational bias--variance trade for BN parents~\cite{Assaf2006,Onisko2001}.

Capital-project open data also embed selection into what gets written. Reasons may be more complete for politically salient projects or for projects already under scrutiny. If so, factor presence is not missing completely at random. The study cannot fully correct that selection from the released fields; it can only avoid compounding it with leakage and silent row deletion~\cite{Woldesenbet2016,Flyvbjerg2009}.

Relative to marine delay BN work, the present evidence channel is weaker but the sample is broader in project count~\cite{Comu2023}. Relative to elicitation-heavy mass-housing BN work, the present study trades expert depth for observational scale and for a budgeted policy evaluation that those studies typically omit~\cite{Khodabakhshian2025}. Relative to progress-prediction ML, it adds calibration and inspection layers that accuracy-centric bake-offs often skip~\cite{Tagharobi2025,Egwim2021}.

The phrase ``operating envelope'' is used deliberately in place of ``robustness'' or ``generalizability.'' Robustness would require stress tests under defined perturbations; generalizability would require external cohorts. What is demonstrated is the envelope on this dataset under the stated clamps, splits, and budgets. That lexical restraint is part of claim calibration~\cite{Gneiting2007,Memarzadeh2016}.

Results figures are sequenced to match inference rather than to form a gallery: literature yardstick, metric pack, ranking curves, error counts, calibration, net benefit, uncertainty, budgets, detection steps, lead time, mapping quality, and utility~\cite{Saito2015,Vickers2006}.

If a reviewer asks what single experiment most justifies the paper, the answer is the joint TEST evaluation of the champion and the budgeted BN scoreboard under PID-held-out open data. Either piece alone would be a smaller contribution; together they answer complementary RQs and reveal the complementary success/boundary pattern that defines the study~\cite{Luu2009,Howard1966,Egwim2021}.

The bibliography emphasizes peer-reviewed journal sources in project management, reliability, forecasting, and construction informatics. Conference-only machine-learning citations were avoided as primary evidence for the journal-count requirement, even when algorithms are standard. Friedman’s gradient-boosting paper remains an appropriate journal foundation for the champion family~\cite{Friedman2001}.

Taken as a whole, the Results and Discussion argue that experimentally measured boundaries are first-class scientific outputs. The VoI tie is not a failed subplot; it is a central finding about inspectable evidence sparsity. The calibrated early-warning success is not a consolation prize; it is the actionable layer agencies can use while evidence regimes remain sparse~\cite{Song2021,Gneiting2007,Tagharobi2025}.

One additional quantitative nuance concerns Brier movement under budget. Although classification metrics are flat across $B$, Brier improves from approximately $0.102$ at $B{=}1$ to approximately $0.097$ at $B{\geq}2$ for the consensus and fixed policies. That movement shows that extra inspections still refine probability quality even when they do not flip alerts. Proper scoring therefore detects a benefit that accuracy-only scoreboards would miss~\cite{Brier1950,Murphy1973,Gneiting2007}.

Disagreement VoI's slightly lower Brier at $B{=}1$ ($0.100$ versus $0.102$) and slightly higher AUC ($0.873$ versus $0.872$) are consistent with a mild prior-induced ranking effect that does not survive into alert counts. Reporting these third-decimal differences without elevating them to policy recommendations is intentional: they are real in the replay logs, but they are not operationally decisive at $\tau{=}0.51$~\cite{Song2021,Memarzadeh2016}.

The early-warning log-loss of $0.146$ complements Brier by penalizing confident misclassification more heavily. Together with ECE $0.030$, the probabilistic triad indicates that the champion's probabilities are usable for triage under the reported resolution. Agencies that need finer reliability at the extreme upper tail should monitor bin-wise calibration as portfolios evolve~\cite{GneitingRaftery2007,Vickers2006}.

Finally, the study's experimental template is intended to be reusable. A new city with a comparable multi-period schedule history can adopt the same PID split, factor ontology (or a documented adaptation), Noisy-OR parameterization, budgeted policy set, and evaluation pack. If that replication yields VoI differentiation, the NYC tie will be understood as regime-specific; if it again yields a tie, sparse administrative reasons will look like a broader boundary condition for sequential delay-factor inspection~\cite{Woldesenbet2016,Luu2009,Regona2024}.

A concluding interpretive synthesis ties the mathematics back to the RQs without introducing new numbers. RQ1 asked whether VoI-guided inspection improves budgeted detection under equal $B$; the scoreboard answers no under absence-clamped NYC reasons, with the mechanistic explanation that inspectable sets are too small for order to matter. RQ2 asked whether disagreement priors change outcomes; alert metrics say no at $\tau{=}0.51$, while third-decimal AUC/Brier shifts remain below operational relevance. RQ3 asked whether prefix early warning can be literature-competitive and calibrated; the champion metrics, curves, and bootstrap intervals answer yes on this dataset under a scoped comparison. The joint reading is that calibrated triage is currently the high-value layer, while budgeted BN inspection is the auditable layer whose sequential intelligence becomes identifiable only when evidence menus thicken~\cite{Howard1966,Luu2009,Egwim2021,Song2021,Gneiting2007,Tagharobi2025}.

That reading also explains why the abstract allocates space to both the positive early-warning findings and the negative VoI-comparative findings. An abstract that reported only accuracy gains would misrepresent the experimental campaign. An abstract that reported only the VoI tie would understate the actionable triage result. Balanced reporting of both is required for claim calibration and for useful guidance to agencies that must allocate attention under incomplete administrative evidence~\cite{Vickers2006,Memarzadeh2016,Regona2024}.
